\documentclass[instructions.tex]{subfiles}
\begin{document}
 
\chapter{On ne doit pas craindre la confession, ni s’en éloigner.}
 
Qu’est-ce qui vous éloigne de la confession? est-ce la peine de dire vos péchés? Mais pourquoi regardez-vous comme une peine, une chose qui fait la consolation des autres, qui doit vous décharger d’un poids énorme, et vous réconcilier à Dieu?
 
Est-ce la crainte de faire une mauvaise confession, et de n’être pas disposé? Mais plus vous différerez, moins vous serez disposé; parce que, plus on s’éloigne de Dieu, plus il est difficile de s’en rapprocher : d’ailleurs, on n’apprend pas à bien faire une chose, en la faisant rarement.
 
Est-ce parce que vous ne voulez pas encore quitter vos péchés? Ah! malheureux, vous voulez donc vivre dans l’état de damnation, et dormir tranquillement sur le bord de l’abîme? Vous ne vous contentez pas d’avoir donné à votre âme le coup de la mort; vous voulez la laisser pourrir dans ses cicatrices. Oh! que vous êtes à plaindre, et d’autant plus à plaindre, que vous ne vous plaignez pas vous-même!
 
Est-ce parce que vous avez honte de déclarer vos péchés? mais pourquoi auriez-vous honte de confesser à un homme ce que vous n’avez pas eu honte de faire sous les yeux de Dieu? Si vous deviez déclarer vos péchés à un ange; mais c’est à un homme faible comme vous, qui connaît les misères du cœur humain, qui doit vous entendre avec des sentimens de charité, et sous le secret le plus inviolable.
 
Saint Paul et saint Augustin ont fait connaître à tout l’univers, dans leurs écrits, les égaremens de leur vie, et les ont publiés en répandant des torrens de larmes. Pourquoi auriez-vous honte de confesser dans le secret des péchés qui sont peut-être moins énormes que les leurs? Ah! vous n’avez pas éprouvé combien il est doux de déclarer ses crimes pour décharger sa conscience d’un poids qui l’accable, et pour mettre son âme en paix.
 
Vous craignez peut-être que le confesseur ne soit étonné d’entendre vos péchés; détrompez-vous. Les médecins n’ont point d’horreur de voir des plaies et des ulcères, de même les confesseurs ne sont point surpris de voir les plaies des consciences; plus une conscience est ulcérée, plus aussi ont-ils de joie de la guérir.
 
Tous les sujets de crainte qui vous éloignent du Sacrement ne sont que dans votre imagination. Tout est effacé aux yeux de Dieu par le sang de Jésus-Christ; dans la confession, tout est de même effacé aux yeux du confesseur : quand vous lui ouvrez votre cœur, il vous estime, il a du respect pour vous, parce qu’il voit en vous les effets de la miséricorde de Dieu; il est rempli de consolation, voyant que Dieu se sert de son ministère pour votre salut; il a pour vous la tendresse d’un père, parce qu’il vous fait renaître à la grâce, et que vous avez pour lui une confiance qu’on n’aurait pas pour son propre père.
 
Si vous ne pouvez vaincre la répugnance que vous avez de dire vos péchés à votre confesseur ordinaire (ce qui serait peut-être le mieux), vous pouvez vous adresser à un autre, pourvu que vous ayez un propos sincère de charger de vie. Il vaut mieux changer de confesseur, que de faire une confession sacrilége. Mais en changeant de confesseur, n’en prenez pas occasion de pécher avec plus de liberté.
 
\section{Résolutions}
 
\primo Non, ni la crainte d’accuser mes péchés, ni la honte qui en accompagne la déclaration, ni la considération de ce que pourra penser de moi le confesseur, ne m’empêcheront jamais de mettre ordre à ma conscience par une bonne confession.
 
\secundo Et je regarderai toujours le ministre du sacrement comme le lieutenant de Dieu, qui me bénit, qui m’écoute, m’interroge, m’exhorte, m’impose une satisfaction, et m’absout en son nom. 

\prayer{Je vous remercie, ô mon Dieu! de ce que vous nous avez donné sur la terre des guides qui y tiennent votre place et que vous avez revêtus de vos pouvoirs, pour nous réconcilier avec vous.}
 
\end{document}
